% sec10_numerics.tex
% Section 10: numerical illustrations.
% Labels defined here: sec:numerics, tab:numerics, tab:onen.

\section{Numerical illustrations}\label{sec:numerics}

This section compares the theorems with computed recurrence coefficients. Table~\ref{tab:numerics} shows the normalised
deviations controlled by Theorem~\ref{thm:analytic} and Corollary~\ref{cor:pure} for a factor $\ell$ that grows faster
than every power of $\log k$, and Table~\ref{tab:onen} compares fitted $1/n$ coefficients with Theorem~\ref{thm:onen}
for $p>1$ and with Conjecture~\ref{conj:onen} for $p\le1$. The computations are floating-point illustrations; no proof
in this paper uses them.

\emph{Method.} The atoms $k\le K$ alone do not suffice. The spectral weight of $P_n$ near the hard edge sits at indices
$k\asymp n^{1+2/p}$, about $10^{16}$ for $p=\frac12$ and $n=1600$, and a truncated measure has a soft edge at its
smallest atom and obeys a different law. For $p=\frac12$, $s=\frac32$, the factor $\ell$ of Table~\ref{tab:numerics} and
the atoms $k\le2\cdot10^5$ alone, $n(n^pb_n/c(p)-1)$ equals $-1.81$, $-13.3$ and $-105.9$ at $n=100$, $400$ and $1600$.
We therefore replaced the atoms $k>K$ by a quadrature of the continuous tail $\int_{K+1/2}^{\infty}t^{-s}\ell(t)\,
\delta_{t^{-p}}\diff t$ and computed the Jacobi matrix of the resulting discrete measure by the updating algorithm that
adds one atom at a time and restores tridiagonal form by plane rotations, in extended precision (the precision and the
truncation checks are given with each table). The scripts and their outputs are available as described in the statement
on data and code availability at the end of the paper.

\subsection{\texorpdfstring{The rate $O(1/n)$}{The rate O(1/n)}}
Table~\ref{tab:numerics} compares the factor $\ell(k)=\exp\sqrt{\log(k+e)}$ with $\ell\equiv1$ for $p=\frac12$, $1$, $2$
and $s=2$. Here $\ell(k)=\exp\Lambda(\log k)$ with $\Lambda(\xi)=\sqrt{\log(e^\xi+e)}$, whose singularities lie on the
lines $|\im\xi|=\pi$, so $\ell$ is analytic of angle $\vartheta$ for every $\vartheta<\pi$
(Definition~\ref{def:analytic}). By Theorem~\ref{thm:analytic} and Corollary~\ref{cor:pure} the normalised deviations
$n(n^pb_n/c(p)-1)$ and $n(a_n/(2b_n)-1)$ stay bounded. The Bessel orders $\alpha=(s-1-p)/p$ are $\alpha=1$, which lies
outside the window of Theorem~\ref{thm:window} and is reached through the shift (Theorem~\ref{thm:shift}) from
$s=\frac32$; $\alpha=0$, the logarithmic case of the Bessel model; and $\alpha=-\frac12$.

The tail was integrated by Gauss--Legendre quadrature on logarithmic panels (ratio $1.02$, eight nodes per panel) up to
$T=10^{48}$ for $p=\frac12$ and $T=10^{16}$ for $p=1,2$, and the recurrence was run in long double precision. We
repeated each computation with $K=10^5$, with $K=2\cdot10^5$, and with $K=2\cdot10^5$ and $T$ reduced by a factor
$10^4$. The largest change of any tabulated entry under these variations is $9.5\cdot10^{-8}$ for $p=\frac12$,
$1.9\cdot10^{-8}$ for $p=1$ and $1.7\cdot10^{-11}$ for $p=2$. The column $\ell\equiv1$ at $p=2$ agrees in every digit
shown with the closed form of Remark~\ref{rem:lommel}.

\begin{table}[ht]
\centering
\small
\begin{tabular}{rrrrrr}
\toprule
 & \multicolumn{2}{c}{$n\bigl(n^pb_n/c(p)-1\bigr)$} & \multicolumn{2}{c}{$n\bigl(a_n/(2b_n)-1\bigr)$} & \\
\cmidrule(lr){2-3}\cmidrule(lr){4-5}
$n$ & $\ell$ & $\ell\equiv1$ & $\ell$ & $\ell\equiv1$ & $\frac12\Lambda'(\log n)$\\
\midrule
\multicolumn{6}{l}{$p=\frac12$, $s=2$, $\alpha=1$ (reached through the shift)}\\
 100 & $-0.4039$ & $-0.5069$ & $-0.2442$ & $-0.2408$ & $0.1131$\\
 200 & $-0.4069$ & $-0.5052$ & $-0.2470$ & $-0.2453$ & $0.1070$\\
 400 & $-0.4101$ & $-0.5039$ & $-0.2484$ & $-0.2477$ & $0.1014$\\
 800 & $-0.4133$ & $-0.5028$ & $-0.2492$ & $-0.2488$ & $0.0963$\\
1600 & $-0.4163$ & $-0.5020$ & $-0.2496$ & $-0.2494$ & $0.0919$\\
\midrule
\multicolumn{6}{l}{$p=1$, $s=2$, $\alpha=0$}\\
  50 & $-0.3821$ & $-0.4934$ & $-0.4954$ & $-0.4950$ & $0.1191$\\
 100 & $-0.3893$ & $-0.4967$ & $-0.4978$ & $-0.4975$ & $0.1131$\\
 200 & $-0.3958$ & $-0.4983$ & $-0.4989$ & $-0.4987$ & $0.1070$\\
 400 & $-0.4015$ & $-0.4992$ & $-0.4994$ & $-0.4994$ & $0.1014$\\
 800 & $-0.4065$ & $-0.4996$ & $-0.4997$ & $-0.4997$ & $0.0963$\\
\midrule
\multicolumn{6}{l}{$p=2$, $s=2$, $\alpha=-\frac12$}\\
  50 & $-0.3794$ & $-0.4938$ & $-0.9756$ & $-0.9780$ & $0.1191$\\
 100 & $-0.3868$ & $-0.4969$ & $-0.9878$ & $-0.9889$ & $0.1131$\\
 200 & $-0.3936$ & $-0.4984$ & $-0.9939$ & $-0.9944$ & $0.1070$\\
 400 & $-0.3995$ & $-0.4992$ & $-0.9970$ & $-0.9972$ & $0.1014$\\
 800 & $-0.4046$ & $-0.4996$ & $-0.9985$ & $-0.9986$ & $0.0963$\\
\bottomrule
\end{tabular}
\caption{Normalised deviations from the pointwise law for $\sigma_{p,2,\ell}$, with $\ell(k)=e^{\sqrt{\log(k+e)}}$, and
for $\ell\equiv1$, computed in extended precision from the atoms $k\le K$ and a quadrature of the continuous tail. All
entries are stable under doubling of $K$ and under a change of the tail cut-off (see the text). Boundedness is what
Theorem~\ref{thm:analytic} and Corollary~\ref{cor:pure} assert. The limits $-(s-1)/2$ and $-p/2$ in the columns
$\ell\equiv1$ are proved for $p=2$ (Theorem~\ref{thm:onen}) and observed, not proved, for $p=\frac12$ and $p=1$. The
slow drift of the first column is the effect of $\Lambda'(\log n)$ (Conjecture~\ref{conj:tilt}).}
\label{tab:numerics}
\end{table}

All entries stay bounded, as the theorems require. For the pure powers the two normalised deviations approach
$-(s-1)/2=-\frac12$ and $-p/2$ in all three geometries. At $p=1$ and $p=2$ the distance of the column $\ell\equiv1$ for
$b_n$ from $-\frac12$ roughly halves when $n$ doubles, as a remainder $O(n^{-2})$ in $n^pb_n/c(p)$ requires; at
$p=\frac12$ it shrinks only by a factor close to $2^{-1/2}$, which suggests a remainder of order $n^{-1-p}$ there. For
$\ell$ the deviation of $n^pb_n$ is smaller in absolute value, and its difference from the pure power drifts slowly
towards $0$, roughly like $\frac12\Lambda'(\log n)$ (last column); see Conjecture~\ref{conj:tilt}. The factor $\ell$ has
almost no effect on $a_n/(2b_n)$ at order $1/n$.

\emph{A factor analytic only in a thin strip.} Theorem~\ref{thm:analytic} needs $\ell$ to be analytic of some angle
$\vartheta>0$ only. As a test we took $p=2$ and
\[
  \Lambda(\xi)=\sqrt{\xi+1}+(\xi+2)^{-1/2}\log(1-q\cos2\xi),\qquad q=1/\cosh(0.6).
\]
Here $1-q\cos2\xi$ vanishes at $\xi=\pi m\pm0.3i$ for every integer $m$, so $\Lambda$ is analytic of angle $\vartheta$
exactly for $\vartheta<0.3$, and the singularities are not confined to a bounded part of the strip. The factor
$\ell(k)=e^{\Lambda(\log k)}$ grows like $e^{\sqrt{\log k}}$ and carries a log-periodic oscillation of relative size
$(\log k)^{-1/2}$. With the same numerical scheme, at $n=50\cdot2^{j/4}$, $0\le j\le24$ (so $50\le n\le3200$), the
deviation $n(n^2b_n/c(2)-1)$ oscillates log-periodically between $-0.6031$ and $-0.0293$ for $s=2$
($\alpha=-\frac12$) and between $-1.1014$ and $-0.5494$ for $s=3$ ($\alpha=0$). Its maximum modulus is $0.598$ on
$50\le n<400$ and $0.589$ on $400<n\le3200$ for $s=2$, and $1.097$, respectively $1.088$, for $s=3$, while
$n(a_n/(2b_n)-1)$ stays between $-1.0091$ and $-0.9505$. The range of $n$ covers only about $1.3$ periods of the
oscillation (period $\pi$ in $\log n$), so this is weak evidence against growth. The deviations are bounded, as
Theorem~\ref{thm:analytic} asserts, and the oscillation of the $1/n$ term reflects that of $\Lambda'$, which here has no
sign.

\subsection{\texorpdfstring{The $1/n$ coefficients}{The 1/n coefficients}}
Write $n^pb_n/c(p)=1+d/n+\cdots$ and $a_n/(2b_n)=1+e/n+\cdots$ for the pure powers $\sigma_{p,s}$.
Theorem~\ref{thm:onen} gives $d=-(s-1)/2$ and $e=-p/2$ for $p>1$, and Conjecture~\ref{conj:onen} asserts the same for
$p\le1$. Table~\ref{tab:onen} lists $d$ and $e$ fitted to the computed coefficients with $n\le1600$. The tail was
integrated by a Gauss--Jacobi rule in the variable $x=t^{-p}$ that is exact on all polynomials of the degree needed, so
that the remaining error comes from replacing the atoms $k>K$ by the integral; the recurrence was run in quadruple
precision, with $K=1.6\cdot10^5$ for $p\ge1$ and $K=3.2\cdot10^5$ for $p=\frac12$, and repeated with $2K$. At
$(p,s)=(2,2)$ the computed $b_n$ and $a_n$ agree with the closed form of Remark~\ref{rem:lommel} for $n\le400$ to a
relative error $6.6\cdot10^{-15}$ (at $K=4\cdot10^4$). The fits are least-squares fits in the powers
$n^{-1-i-jp}$, $i,j\ge0$, with several numbers of terms and fitting windows; the table gives the median of the estimates
and half their range as an error estimate. These estimates are not bounds.

\begin{table}[ht]
\centering
\small
\begin{tabular}{rrrrrrrr}
\toprule
 & & \multicolumn{3}{c}{coefficient $d$ of $n^pb_n/c(p)$} & \multicolumn{3}{c}{coefficient $e$ of $a_n/(2b_n)$}\\
\cmidrule(lr){3-5}\cmidrule(lr){6-8}
$p$ & $s$ & fitted & estimate & $-\frac{s-1}2$ & fitted & estimate & $-\frac p2$\\
\midrule
\multicolumn{8}{l}{$p>1$: Theorem~\ref{thm:onen}}\\
$\frac32$ & $\frac32$ & $-0.250000000$ & $3.4\cdot10^{-9}$ & $-0.25$ & $-0.750000000$ & $1.9\cdot10^{-8}$ & $-0.75$\\
$\frac32$ & $2$ & $-0.500000000$ & $4.8\cdot10^{-9}$ & $-0.5$ & $-0.750000000$ & $6.9\cdot10^{-9}$ & $-0.75$\\
$\frac32$ & $3$ & $-1.000000000$ & $1.6\cdot10^{-8}$ & $-1$ & $-0.750000000$ & $1.4\cdot10^{-8}$ & $-0.75$\\
$\frac32$ & $4$ & $-1.500000000$ & $1.4\cdot10^{-8}$ & $-1.5$ & $-0.750000000$ & $1.6\cdot10^{-7}$ & $-0.75$\\
$2$ & $\frac32$ & $-0.250000000$ & $6.0\cdot10^{-10}$ & $-0.25$ & $-1.000000000$ & $1.6\cdot10^{-9}$ & $-1$\\
$2$ & $2$ & $-0.500000000$ & $3.4\cdot10^{-11}$ & $-0.5$ & $-1.000000000$ & $1.1\cdot10^{-10}$ & $-1$\\
$2$ & $3$ & $-1.000000000$ & $3.2\cdot10^{-11}$ & $-1$ & $-1.000000000$ & $1.2\cdot10^{-11}$ & $-1$\\
$2$ & $4$ & $-1.500000000$ & $8.6\cdot10^{-11}$ & $-1.5$ & $-1.000000000$ & $2.6\cdot10^{-10}$ & $-1$\\
$3$ & $\frac32$ & $-0.250000000$ & $2.9\cdot10^{-12}$ & $-0.25$ & $-1.500000000$ & $4.5\cdot10^{-10}$ & $-1.5$\\
$3$ & $2$ & $-0.500000000$ & $3.9\cdot10^{-12}$ & $-0.5$ & $-1.500000000$ & $3.2\cdot10^{-10}$ & $-1.5$\\
$3$ & $3$ & $-1.000000000$ & $2.2\cdot10^{-11}$ & $-1$ & $-1.500000000$ & $1.0\cdot10^{-10}$ & $-1.5$\\
$3$ & $4$ & $-1.500000000$ & $8.1\cdot10^{-11}$ & $-1.5$ & $-1.500000000$ & $1.7\cdot10^{-11}$ & $-1.5$\\
$4$ & $\frac32$ & $-0.250000000$ & $1.8\cdot10^{-12}$ & $-0.25$ & $-2.000000000$ & $1.1\cdot10^{-9}$ & $-2$\\
$4$ & $2$ & $-0.500000000$ & $1.5\cdot10^{-12}$ & $-0.5$ & $-2.000000000$ & $8.8\cdot10^{-10}$ & $-2$\\
$4$ & $3$ & $-1.000000000$ & $1.3\cdot10^{-11}$ & $-1$ & $-2.000000000$ & $4.6\cdot10^{-10}$ & $-2$\\
$4$ & $4$ & $-1.500000000$ & $5.7\cdot10^{-11}$ & $-1.5$ & $-2.000000000$ & $1.7\cdot10^{-10}$ & $-2$\\
\midrule
\multicolumn{8}{l}{$p\le1$: support for Conjecture~\ref{conj:onen}}\\
$\frac12$ & $\frac32$ & $-0.249999998$ & $3.8\cdot10^{-7}$ & $-0.25$ & $-0.250000021$ & $2.7\cdot10^{-7}$ & $-0.25$\\
$\frac12$ & $2$ & $-0.500000014$ & $5.7\cdot10^{-7}$ & $-0.5$ & $-0.249999950$ & $6.3\cdot10^{-6}$ & $-0.25$\\
$\frac12$ & $3$ & $-1.000000247$ & $2.9\cdot10^{-5}$ & $-1$ & $-0.249999294$ & $9.3\cdot10^{-5}$ & $-0.25$\\
$\frac12$ & $4$ & $-1.500001361$ & $1.3\cdot10^{-4}$ & $-1.5$ & $-0.249996155$ & $3.4\cdot10^{-4}$ & $-0.25$\\
$1$ & $\frac32$ & $-0.250000000$ & $1.3\cdot10^{-10}$ & $-0.25$ & $-0.500000000$ & $3.7\cdot10^{-10}$ & $-0.5$\\
$1$ & $2$ & $-0.500000000$ & $5.6\cdot10^{-10}$ & $-0.5$ & $-0.500000000$ & $1.1\cdot10^{-9}$ & $-0.5$\\
$1$ & $3$ & $-1.000000000$ & $5.6\cdot10^{-10}$ & $-1$ & $-0.500000000$ & $1.1\cdot10^{-9}$ & $-0.5$\\
$1$ & $4$ & $-1.500000000$ & $2.1\cdot10^{-9}$ & $-1.5$ & $-0.500000000$ & $8.5\cdot10^{-9}$ & $-0.5$\\
\bottomrule
\end{tabular}
\caption{Fitted $1/n$ coefficients $d$ and $e$ for the pure powers $\sigma_{p,s}$, from the coefficients with
$n\le1600$, against the values of Theorem~\ref{thm:onen} ($p>1$) and of Conjecture~\ref{conj:onen} ($p\le1$). The
estimate is half the range of the fits; the change under doubling of the truncation index $K$ is of the same order or
smaller. Neither is a bound.}
\label{tab:onen}
\end{table}

For $p>1$ every fitted value agrees with Theorem~\ref{thm:onen} in all nine digits shown. For $p=1$ the same holds, and
for $p=\frac12$ the discrepancy is within the error estimate; at $p=\frac12$ the fits need the non-integer powers
$n^{-1-jp}$, in line with the slower convergence seen in Table~\ref{tab:numerics}. At $(p,s)=(2,2)$ the exact values
$d=-\frac12$ and $e=-1$ also follow from the closed form of Remark~\ref{rem:lommel}. That closed form was checked in
interval (ball) arithmetic for $1\le n<300$: the coefficients computed by the Chebyshev algorithm from the moments
$\zeta(2j+2)$ have rigorous enclosures that agree with it to a relative accuracy better than $10^{-2000}$.
